\unnumberedChapter{Hardware Extension Module Design}

// ??? rework... we have node and channels which are the path to hardware boards.


\unnumberedSection{Requirements}


// ??? repeat picture from general concepts part ...



Because of the large variety of extension boards a software layer is needed that unifies how the individual access points such as a digital input or output is accessed by the LCS node firmware. There should be an easy and uniform way to address these endpoints, i.e. there should be some kind of an \textbf{extension library} interface common to all extension boards designed.


\unnumberedSection{Concepts}

// ??? nodes, ports and channels...

To accommodate the requirements, each extension board has to a hardware mechanism that uniquely identifies the board and its capabilities. It is a not a requirement that such a board has processing power or retains any state. This is the responsibility of the controller board. When power is applied to the main controller the reset or restart sequence will attempt to discover all extension boards connected. For this to work, each board has a non-volatile memory structure that describes the board capabilities. And this memory needs to be at a known place, i.e. I2C address, such that this information can be read without knowing any further details. For this purpose, extension boards will have a NVM chip, that contains the board description.

The NVM chip is "programmed" by writing the relevant data to the chip. A write-enable jumper is used to allow writing and then disable all further changes during operation. There is no need for a special programmer, any main controller board can do this via the I2C lines of the extension connector. With using the NV chip, there is no need to have any further DIP switches or alike on the board. In fact, a memory chip allows for even more precise options to configure. And it is cheaper that even a 4-DIP switch.

After programming the NVM chip will contains all information needed to address the extension board functions. Typically, one or more I2C addressable chips are on the board and the NVM tells us what their address is. Combining the chip foxed I2C address portion, the board address portion and the individual chip select portion, an I2C bus unique address is formed. 

The software library can now access a chip. However, chips vary greatly in their functions and how they are controlled by software. A concept is needed to communicate with these chops at  higher abstraction level. Think of a driver in operating systems. A driver offers a set of routines, such as open, read, write and close and takes care of the underlying details how to talk to the physical object. We will go a similar route. There is a high level library, the extension library that offers a high level view of extension board capabilities. Upon node reset and restart, the controller board will query the NVM chips on the extension bod and dynamically configure the driver code for this board.

In addition to the board address a large variety of inputs, outputs and perhaps functions need to be address too. We will call them endpoints. Consider  16-port digital IO chip, such as the MCP23017. It will have 16 IO pins, which we call 16 endpoints. Uniquely identifying an and point is to know the board ID and the endpoint Id on that board. Again, all this information is stored in the NVM  chip at extension board conjuration time. An endpoint could be a variety of things. For example a plan digital output, a servo output, and analog input, and so on. An endpoint could also be a logical entity that groups multiple pins or issue a series of commands to the extension board.

For a new board design, the development process is therefore to define the NVM data that represent the hardware capabilities  and to provide a driver for the board capabilities. From thereon, the firmware designer can use the extension boards with simple but powerful commands. How does this all map to LCS nodes and ports ? Well, the firmware designer has to provide the code that makes the respective driver calls when receiving node and port control commands and events.


\unnumberedSection{Multiple Extension Boards}

// ??? they are encoded via the port and channels number. One can have boards directly connected to the extension connector and via a cable as satellites...


\unnumberedSection{Board Discovery and Setup}

// ??? we do not need a discovery... channels do the trick ... 
// ??? we try to read from a channels and if response -> found one .



\unnumberedSection{Summary}

LCS nodes consist of a main controller board and extension boards. That concept can be found throughout all what is presented in this book. This chapter gave an insight how many different extension boards are connected and managed. We are now read to look at actual extension board implementations in the chapters to come.


??? a newer concept, started with "basic" extension, but this is more powerful.



